% 大爆炸（综述）
% license CCBYSA3
% type Wiki

本文根据 CC-BY-SA 协议转载翻译自维基百科\href{https://en.wikipedia.org/wiki/Big_Bang}{相关文章}。

\begin{figure}[ht]
\centering
\includegraphics[width=8cm]{./figures/7a2c26b638cb807a.png}
\caption{宇宙膨胀的时间线：其中空间在每个阶段以圆形剖面示意表示。左侧展示了宇宙暴胀时期的剧烈膨胀；中间部分则表示宇宙膨胀的加速阶段（艺术示意图；时间与尺度均非按比例绘制）。} \label{fig_BBang_1}
\end{figure}
“大爆炸”（The Big Bang）是一种物理理论，用以描述宇宙如何从最初高密度、高温的状态中膨胀而来。基于大爆炸概念的各种宇宙学模型能够解释广泛的天文现象\(^\text{[1][2][3]}\)包括轻元素的丰度、宇宙微波背景辐射（CMB），以及宇宙的大尺度结构。宇宙的均匀性——即所谓的视界与平坦性问题——通过“宇宙暴胀”得到解释：那是一段发生在宇宙最早期的加速膨胀阶段。对宇宙膨胀速率的精确测量将最初奇点的时间推算为约 137.87 ± 0.02 亿年前，这被认为是宇宙的年龄。大量实证证据强烈支持这一“大爆炸事件”，该理论如今已被广泛接受。\(^\text{[4]}\)

根据已知物理定律，将这种宇宙膨胀过程反向外推，模型显示宇宙早期处于极端炽热且高密度的原始状态。物理学目前仍缺乏一种被广泛接受的理论来完整描述大爆炸最初阶段的条件。\(^\text{[5]}\)随着宇宙膨胀，它逐渐冷却到足以形成亚原子粒子，随后又形成原子。这些原初元素——主要是氢，辅以少量氦和锂——在暗物质的引力作用下聚集，最终形成早期的恒星与星系。对超新星红移的观测表明，宇宙的膨胀正在加速，这一现象被归因于一种被称为“暗能量”的概念。

“宇宙膨胀”的思想最早由物理学家亚历山大·弗里德曼（Alexander Friedmann）于 1922 年提出，他在数学上推导出了“弗里德曼方程”。\(^\text{[6][7][8][9]}\)宇宙膨胀的最早观测依据被称为“哈勃定律”（Hubble’s Law），由物理学家爱德文·哈勃（Edwin Hubble）于 1929 年发表，他发现星系正以与其距离成正比的速率远离地球。独立于弗里德曼的理论工作与哈勃的观测，物理学家乔治·勒梅特（Georges Lemaître）于 1931 年提出，宇宙起源于一个“原初原子”（primeval atom），从而引入了现代意义上的“大爆炸”概念。1964 年，人们发现了宇宙微波背景辐射。随后几年中，对该辐射的观测表明：它在天空各方向上高度均匀，且能量—强度曲线的形状都与“大爆炸模型”预测的早期高温高密度状态相一致。至 20 世纪 60 年代末，大多数宇宙学家已确信竞争性的“稳恒态宇宙”模型是错误的。\(^\text{[10]}\)

然而，观测宇宙中仍存在若干现象尚未能由“大爆炸模型”充分解释。其中包括：物质与反物质数量不对称（即重子不对称性）、星系周围暗物质的具体性质，以及“暗能量”的起源。\(^\text{[11]}\)
\subsection{模型的特征}
\subsubsection{基本假设}
大爆炸宇宙学模型依赖于三项主要假设：物理定律的普适性、宇宙学原理，以及物质可被视作理想流体。\(^\text{[12]}\)物理定律的普适性是相对论的基本原则之一。宇宙学原理指出，在大尺度上，宇宙是均匀且各向同性的——无论从哪个方向观察、位于何处，宇宙都呈现相同的总体结构。\(^\text{[13]}\)所谓“理想流体”，是指没有黏度的流体，其压力与密度成正比。\(^\text{[14]: 49 }\)

这些假设最初被视为公设，但后来科学家努力对其逐一加以验证。例如，第一个假设（物理定律的普适性）已通过观测得到检验：在宇宙大部分存在时期内，精细结构常数（fine-structure constant）的最大可能偏差量级仅为 $10^{-5}$。\(^\text{[15]}\)支撑这些模型的关键物理定律——广义相对论——已在太阳系与双星系统尺度上通过了严格的检验。\(^\text{[16][17]}\)宇宙学原理亦得到高度验证：通过对宇宙微波背景（CMB）温度的观测，其一致性水平达到 $10^{-5}$。截至 1995 年，对 CMB 视界尺度的测量表明，宇宙的非均匀性上限约为 10\%。\(^\text{[18]}\)
\subsubsection{膨胀预测}
宇宙学原理极大地简化了广义相对论方程，得出了用于描述宇宙几何结构的弗里德曼–勒梅特–罗伯逊–沃尔克度规（FLRW metric）；结合理想流体假设，可推导出弗里德曼方程，用于描述宇宙几何随时间的演化。在这一层次的描述中，唯一的自由参数是质量-能量密度：宇宙的几何形态与其膨胀行为均由密度直接决定。\(^\text{[19]: 73 }\)大爆炸宇宙学的所有主要特征都与这些结果密切相关。\(^\text{[14]: 49 }\)
\subsubsection{质量–能量密度}
\subs\begin{figure}[ht]
\centering
\includegraphics[width=8cm]{./figures/818e9d5df373b8f0.png}
\caption{} \label{fig_BBang_2}
\end{figure}
在大爆炸宇宙学中，质量–能量密度决定了宇宙的形状与演化。通过将天文观测结果与已知的热力学和粒子物理定律结合，宇宙学家推算出了宇宙生命期内密度的组成比例。在当前的宇宙中，可见物质（如恒星、行星等）所占的密度不到 5\%；暗物质约占 27\%；而暗能量则占据剩余的 68\%。\(^\text{[20]}\)
\subsubsection{视界}
大爆炸时空的一个重要特征是粒子视界的存在。由于宇宙的年龄有限，且光速也是有限的，因此可能存在一些过去的事件，其光尚未有足够的时间到达地球。这就为我们能观测到的最遥远天体设置了一个限制，即“过去视界”。  
相反，由于空间在不断膨胀，更遥远的天体正以越来越快的速度远离我们，因此我们今天发出的光可能永远无法“追上”那些极远的天体。这就定义了一个“未来视界”，它限制了我们未来能够影响到的事件范围。  
这两种视界的存在都取决于描述宇宙膨胀的弗里德曼–勒梅特–罗伯逊–沃尔克（FLRW）度规的具体形式。\(^\text{[21]}\)

根据我们对早期宇宙的理解，可以推知宇宙存在“过去视界”；但在实践中，我们的观测受到早期宇宙不透明性的限制，因此视野无法在时间上无限向后延伸，尽管该视界在空间上不断退后。若宇宙的膨胀继续加速，还将存在“未来视界”。\(^\text{[21]}\)
\subsubsection{热化过程}
早期宇宙中的某些物理过程相对于宇宙膨胀速率而言过于缓慢，因而无法达到近似热力学平衡；另一些过程则足够迅速，可以实现热化。  
通常用来判断早期宇宙中的某一过程是否达到热平衡的参数，是该过程速率（通常为粒子碰撞频率）与哈勃参数的比值。比值越大，意味着粒子在彼此分离之前有更多时间完成热化过程。\(^\text{[22]}\)

宇宙能量密度各组成部分的相对分布估计图。（2015 年 2 月，欧洲主导的普朗克宇宙学探测计划研究团队发布了新的精确数据：普通物质占 4.9\%，暗物质占 25.9\%，暗能量占 69.1\%。）
\subsection{时间线}
根据大爆炸模型，宇宙在最初阶段极其炽热而致密，此后一直在膨胀与冷却之中。
\subsubsection{奇点}
现有的物理理论无法描述大爆炸发生的确切瞬间。\(^\text{[5]}\)如果仅用经典广义相对论将宇宙的膨胀反推至过去，会得到一个具有无限密度与温度的引力奇点，该奇点存在于有限的过去时间中。\(^\text{[23]}\)然而，这种经典引力理论被认为不足以描述在如此极端条件下的物理。\(^\text{[19]: 275}\) 因此，这个奇点在大爆炸语境中的真实含义仍不明确。\(^\text{[24]}\)

广义相对论所能适用的最早时刻称为普朗克时间（Planck time）。\(^\text{[19]: 274}\) 在此之前，即所谓的普朗克纪元（Planck epoch），宇宙温度接近普朗克温标（约 $10^{32}$ K 或 $10^{28}$ eV），量子引力效应被认为占主导地位。  
迄今为止，人类尚未建立一个被普遍接受的量子引力理论；在高于普朗克能标的区域，未知的物理规律可能会影响宇宙的膨胀历史。
\subsubsection{暴胀与重子生成}
由于缺乏可观测数据，大爆炸最早阶段的研究仍主要基于推测。在最常见的模型中，宇宙在最初阶段具有极高的能量密度、温度与压力，并以极快的速度均匀膨胀与冷却。  
在膨胀约 $10^{-43}$ 秒之前的普朗克纪元中，四种基本相互作用——电磁力、强核力、弱核力与引力——统一为一种力。\(^\text{[25]}\)此时，宇宙的特征尺度长度为普朗克长度（$1.6\times10^{-35}$ 米），温度约为 $10^{32}$ 摄氏度；在这种条件下，“粒子”的概念本身已失去意义。对这一时期的准确理解有待量子引力理论的发展。\(^\text{[26][27]}\)普朗克纪元之后，在约 $10^{-43}$ 秒时进入\emph{大统一纪元}，此时引力从其他三种力中分离出来，因为宇宙温度继续下降。\(^\text{[25]}\)

大约在膨胀 $10^{-37}$ 秒时，一次相变引发了宇宙暴胀（cosmic inflation），期间宇宙以指数级速度膨胀，不受光速不变性的约束，温度下降了约十万倍。这一思想的提出是为了解释平坦性问题（flatness problem）——即宇宙的物质与能量密度非常接近产生平坦几何所需的临界密度。换言之，宇宙整体没有几何曲率。由于海森堡不确定性原理产生的微观量子涨落在暴胀期间被“冻结”，并被放大为后来形成宇宙大尺度结构的种子。\(^\text{[28]}\)在约 $10^{-36}$ 秒时进入电弱纪元（electroweak epoch），此时强核力从其他力中分离，只剩电磁力与弱核力保持统一。\(^\text{[29]}\)

当前可观测到的所有星系的质量–能量最初都集中在一个半径约 $4\times10^{-29}$ 米的球体中，到暴胀结束时该球体的半径约增至 0.9 米。\(^\text{[30]: 202}\) 随后进入再加热阶段（reheating），暴胀子场（inflaton field）衰变，使宇宙温度升高到足以产生夸克–胶子等离子体及各种基本粒子的水平。\(^\text{[31][32]}\)当时温度极高，粒子的随机运动达到相对论速度，各类粒子与反粒子对在碰撞中不断产生与湮灭。\(^\text{[33]}\)在某个阶段，发生了一种未知的反应称为重子生成（baryogenesis），它违反了重子数守恒定律，使夸克与轻子相对于反夸克与反轻子产生极微小的不对称，约为三千万分之一。这种微小的不对称导致了当今宇宙中物质相对于反物质的主导地位。\(^\text{[34]}\)
\subsubsection{冷却}
\begin{figure}[ht]
\centering
\includegraphics[width=8cm]{./figures/c016fce83cce5d6b.png}
\caption{整个近红外天空的全景图展示了银河系之外星系的分布。星系按照红移（redshift）进行了颜色编码。} \label{fig_BBang_3}
\end{figure}
宇宙的密度继续下降，温度不断降低，因此每个粒子的典型能量也随之减小。  
随着对称破缺相变（symmetry-breaking phase transitions）的发生，物理学的基本相互作用以及基本粒子的参数逐渐形成今天的状态；在约 $10^{-12}$ 秒时，电磁力与弱核力相互分离。\(^\text{[29][35]}\)

在约 $10^{-11}$ 秒之后，情况变得不再完全推测性，因为此时的粒子能量已降至可在粒子加速器中达到的范围。大约在 $10^{-6}$ 秒时，夸克与胶子结合形成重子（baryons），如质子与中子。夸克相对于反夸克的微小过剩导致重子相对于反重子的微小过剩。此时温度已不足以产生新的质子–反质子对或中子–反中子对，随之发生了大规模的湮灭反应，最终仅剩原始物质粒子中约 $10^8$ 分之一的部分，而其反粒子全部消失。\(^\text{[36]}\)类似的过程在约 1 秒时也发生于电子与正电子之间。在这些湮灭之后，剩下的质子、中子与电子已不再以相对论速度运动，宇宙的能量密度此时由光子主导（中微子的贡献较小）。

在膨胀几分钟后，当温度约为十亿开（$10^9$ K），宇宙中的物质密度与当今地球大气密度相当时，中子与质子结合生成氘核与氦核，这一过程被称为大爆炸核合成（Big Bang nucleosynthesis, BBN）。\(^\text{[37]}\)大多数质子未被结合，保留为氢核。\(^\text{[38]}\)

随着宇宙继续冷却，在约 5 万年左右，物质的静能密度开始在引力上压倒光子与中微子辐射的能量密度。在约 38 万年时，宇宙冷却至电子与原子核可结合成中性原子（主要为氢）的温度，这一事件称为复合纪元（recombination）。该过程使原本不透明的宇宙变得透明，而在此时期最后一次散射的光子构成了宇宙微波背景辐射（Cosmic Microwave Background, CMB）。\(^\text{[38]}\)
\subsubsection{结构形成}
\begin{figure}[ht]
\centering
\includegraphics[width=8cm]{./figures/a1ebb27bafba8b13.png}
\caption{Abell 2744 星系团——哈勃前沿视野（Hubble Frontier Fields）观测图\(^\text{[39]}\)} \label{fig_BBang_4}
\end{figure}
在复合纪元（recombination epoch）之后，均匀分布的物质中稍微致密的区域通过引力吸引周围物质，从而变得更加致密，逐渐形成气体云、恒星、星系以及今天可观测到的其他天体结构。\(^\text{[33]}\)这一过程的细节取决于宇宙中物质的数量与类型。物质被分为四种可能的类型：冷暗物质（cold dark matter, CDM）、温暗物质（warm dark matter）、热暗物质（hot dark matter）以及重子物质（baryonic matter）。根据威尔金森微波各向异性探测器（Wilkinson Microwave Anisotropy Probe, WMAP）的最佳测量结果，数据与假设暗物质为冷态的 $\Lambda$CDM 模型高度吻合。估计这种冷暗物质约占宇宙总物质/能量的 23\%，而重子物质约占 4.6\%。\(^\text{[40]}\)
\subsubsection{宇宙加速膨胀}
来自 Ia 型超新星与宇宙微波背景辐射（CMB）的独立观测证据表明，当今宇宙由一种被称为暗能量（dark energy）的神秘能量形式主导，这种能量似乎均匀地充满整个空间。观测显示，目前宇宙总能量密度的约 73\% 属于这种形式。在宇宙早期阶段，暗能量可能已经存在，但由于物质彼此更为接近，引力占主导地位，抑制了膨胀。最终，经过数十亿年的膨胀，物质密度相对于暗能量密度的下降，使得宇宙的膨胀开始加速。\(^\text{[11]}\)

在最简单的表述中，暗能量可通过爱因斯坦广义相对论场方程中的宇宙常数项（cosmological constant term）建模，但其组成与机制仍然未知。更一般地说，其状态方程（equation of state）的细节以及与粒子物理标准模型的关系，仍通过观测与理论研究不断探索。\(^\text{[11]}\)

暴胀纪元之后的所有宇宙演化，都可以由宇宙学的ΛCDM 模型（lambda-CDM model）严格描述与建模。该模型结合了量子力学与广义相对论的独立框架。然而，在约 $10^{-15}$ 秒之前的宇宙状态，尚无可直接验证的模型能够加以描述。\(^\text{[41]}\)理解宇宙历史中这最早的阶段，是现代物理学中最伟大的未解难题之一。
\subsection{概念史}
\subsubsection{词源}
英国天文学家弗雷德·霍伊尔（Fred Hoyle）被认为是“Big Bang（大爆炸）”一词的创始者。他于 1949 年 3 月在 BBC 电台的演讲中首次使用这一说法，\(^\text{[42]}\)并说道：“这些理论基于这样一个假设：宇宙中的所有物质都是在遥远过去某一特定时刻，通过一次巨大的爆炸（one big bang）被创造出来的。”\(^\text{[43][44]}\)然而，这个名称直到 1970 年代才真正流行起来。\(^\text{[44]}\)

人们普遍认为，支持另一种“稳态宇宙模型（steady-state cosmological model）”的霍伊尔，最初是带着贬义使用这一说法的，\(^\text{[45][46][47]}\)但霍伊尔本人明确否认了这一点，称他只是为了突出两种宇宙模型的区别而使用了这个生动的比喻。\(^\text{[48][49][51]}\)赫尔格·克拉格（Helge Kragh）指出，认为该词带有贬义的证据“并不令人信服”，并列举了若干迹象表明它并非贬义用语。\(^\text{[44]}\)

“原初奇点（primordial singularity）”有时也被称为“大爆炸（Big Bang）”，\(^\text{[52]}\)但该术语也可以泛指早期宇宙中炽热而致密的阶段。\(^\text{[53]}\)一些学者认为“Big Bang”这个词是一个误称（misnomer），因为它让人联想到爆炸。\(^\text{[44][54]}\)他们指出，爆炸通常意味着向外部空间的扩张，而“大爆炸”实际上描述的只是宇宙内部内容物的内在膨胀。\(^\text{[55][56]}\)另一个由桑托什·马修（Santhosh Mathew）指出的问题是，“bang（爆炸）”一词暗示声音，但声音并不是该模型的重要特征。\(^\text{[46]}\)然而，试图寻找更合适替代词的努力最终未能成功。\(^\text{[44]}\)

根据蒂莫西·费里斯（Timothy Ferris）的说法：\(^\text{[47][57]}\)  

“‘Big Bang’ 这一术语原本带有嘲讽意味，是由弗雷德·霍伊尔创造的，而它的持久流行恰恰证明了霍伊尔爵士的创造力与机智。  
事实上，这个词在一场国际征名比赛中幸存下来——三位评审（电视科学记者休·唐斯（Hugh Downs）、天文学家卡尔·萨根（Carl Sagan）以及我本人）  
从 41 个国家的 13,099 份投稿中筛选，但没有任何一个足够恰当以取代它。  
最终没有获胜者被宣布。无论你喜欢与否，我们都只能继续使用‘Big Bang’这个词。”
\subsubsection{命名前}
\begin{figure}[ht]
\centering
\includegraphics[width=8cm]{./figures/25efff4c27c87e40.png}
\caption{XDF（哈勃极深场，Hubble eXtreme Deep Field）与月球视直径的大小比较图 —— XDF 位于月球左侧偏下方的一个小方框内。在这极小的视野中，包含了数千个星系，而每个星系又由数十亿颗恒星组成。} \label{fig_BBang_5}
\end{figure}
\begin{figure}[ht]
\centering
\includegraphics[width=8cm]{./figures/4350f947c8088aa3.png}
\caption{XDF（2012）观测图 —— 每一个光点都是一个星系，其中一些星系的年龄高达 132 亿年。\(^\text{[58]}\)据估计，整个宇宙包含约 2000 亿个星系。} \label{fig_BBang_6}
\end{figure}
\begin{figure}[ht]
\centering
\includegraphics[width=8cm]{./figures/8fa86fae373275ae.png}
\caption{XDF 图像展示了前景中的完全成熟星系；距今约 50 至 90 亿年的近成熟星系；以及更远处、年龄超过 90 亿年的原始星系（原星系团），其中炽烈闪耀着年轻的恒星。} \label{fig_BBang_7}
\end{figure}
早期的宇宙学模型源自对宇宙结构的观测与理论推演。1912 年，美国天文学家维斯托·斯里弗（Vesto Slipher）首次测量到“旋涡星云”（即今天所称的螺旋星系）的多普勒位移，并很快发现几乎所有此类星云都在远离地球。他当时并未意识到这一现象的宇宙学意义；事实上，在那个时代，这些星云是否属于银河系之外的“岛宇宙”（island universes）仍是极具争议的问题。\(^\text{[59][60]}\) 

十年后，俄国宇宙学家兼数学家亚历山大·弗里德曼（Alexander Friedmann） 从爱因斯坦场方程中推导出弗里德曼方程（Friedmann equations），表明宇宙可能正在膨胀——这一结果与当时阿尔伯特·爱因斯坦所主张的静态宇宙模型相矛盾。\(^\text{[61][62]}\)

1924 年，美国天文学家埃德温·哈勃（Edwin Hubble）通过测定最近螺旋星云的巨大距离，证明这些系统确实是独立的星系。从同年开始，哈勃利用威尔逊山天文台的 100 英寸（2.5 米）胡克望远镜，painstakingly 建立了一系列距离指示器，为后来的“宇宙距离梯”（cosmic distance ladder）奠定了基础。这使他能够估算那些已由斯里弗测得红移的星系距离。1929 年，哈勃发现了星系距离与退行速度之间的相关性——即如今所称的哈勃定律（Hubble's law）。\(^\text{[63][64]}\)

1927 年，比利时物理学家、天主教神父乔治·勒梅特（Georges Lemaître）独立推导出弗里德曼方程，提出星云的退行是由宇宙膨胀造成的。\(^\text{[65][66]}\)基于宇宙学原理（cosmological principle），他预言了哈勃后来实际观测到的关系。\(^\text{[11]}\)1931 年，勒梅特更进一步指出：如果将宇宙的膨胀追溯到过去，意味着越往过去宇宙越小，直到在某一有限的过去时刻，宇宙中所有质量都集中在一个“原初原子”（primeval atom）中，在那里，时间与空间的织构开始存在。\(^\text{[67]}\)

在 1920 至 1930 年代，几乎所有主要的宇宙学家都倾向于永恒稳态宇宙（steady-state universe）。其中许多人批评大爆炸所暗示的“时间起点”会将宗教观念引入物理学；这一反对意见后来被稳态理论的支持者重复引用。\(^\text{[68]}\)这种印象因膨胀宇宙概念的提出者勒梅特本身是一位天主教神父而进一步加深。\(^\text{[69]}\)阿瑟·爱丁顿（Arthur Eddington）同意亚里士多德的观点，认为宇宙没有时间上的起点，即物质是永恒的；他认为时间的开端是“令人厌恶的”（“repugnant”）。\(^\text{[70][71]}\)但勒梅特不同意这一看法，他写道：  

“如果世界起始于一个单一的量子（single quantum），那么在最初，空间与时间的概念将完全失去意义；只有当这个原初量子被分割为足够数量的量子之后，空间与时间的概念才会获得合理的意义。如果这一设想正确，那么世界的开端略早于空间与时间的开端。”\(^\text{[72]}\) 

在 1930 年代，为解释哈勃的观测，人们还提出了其他一些非标准宇宙学（non-standard cosmologies）：如米尔恩模型（Milne model），\(^\text{[73]}\)振荡宇宙模型（oscillatory universe，最初由弗里德曼提出，后由爱因斯坦与托尔曼倡导），\(^\text{[74]}\)以及弗里茨·兹威基（Fritz Zwicky）的“光疲劳假说”（tired light hypothesis）。\(^\text{[75]}\)  

第二次世界大战后，宇宙学出现了两种截然不同的可能性：一种是弗雷德·霍伊尔的稳态模型（steady-state model），其中宇宙膨胀的同时不断生成新物质，使其在任何时刻看起来都大致相同。\(^\text{[76]}\)另一种是勒梅特的膨胀宇宙理论（expanding universe theory），由乔治·伽莫夫（George Gamow）发展与推广，他以此建立了宇宙化学元素丰度的理论，\(^\text{[77]}\)其同事拉尔夫·阿尔弗（Ralph Alpher）与罗伯特·赫尔曼（Robert Herman）则预测了宇宙微波背景辐射的存在。\(^\text{[78]}\)
\subsubsection{作为一个命名模型}
具有讽刺意味的是，正是霍伊尔（Fred Hoyle）本人创造了后来被用于勒梅特（Georges Lemaître）理论的词汇“Big Bang”（大爆炸）。1949 年 3 月，他在 BBC 电台广播中将其称为 “this big bang idea”（“这个大爆炸的想法”）。\(^\text{[49][44][notes 1]}\)在相当长一段时间里，学界的支持在这两种理论（稳态宇宙与大爆炸宇宙）之间分裂。最终，观测证据——尤其是无线电源计数（radio source counts）——开始倾向于大爆炸模型。1964 年宇宙微波背景辐射（CMB）的发现与验证，确立了大爆炸理论作为宇宙起源与演化的最佳模型。\(^\text{[79]}\)

1968 年与 1970 年，罗杰·彭罗斯（Roger Penrose）、斯蒂芬·霍金（Stephen Hawking）以及乔治·F·R·埃利斯（George F. R. Ellis）发表论文，证明在相对论性的大爆炸模型中，数学奇点是不可避免的初始条件。\(^\text{[80][81]}\)随后，从 1970 年代到 1990 年代，宇宙学家们致力于刻画大爆炸宇宙的特征，并解决其中的若干悬而未决的问题。1981 年，艾伦·古斯（Alan Guth）在理论上取得突破，他提出了“大爆炸模型”中的一个关键修正阶段—— 早期宇宙的“暴胀”（inflation）时期，用以解决若干长期存在的理论难题。\(^\text{[82]}\)

与此同时，在这几十年间，观测宇宙学中引发广泛讨论与分歧的两个问题是：哈勃常数（Hubble Constant）的精确数值\(^\text{[83]}\)，以及宇宙的物质密度（在发现暗能量之前，人们认为它是预测宇宙最终命运的关键参数）。\(^\text{[84]}\)

自 1990 年代末以来，由于望远镜技术的进步与来自多颗卫星（如宇宙背景探测者 COBE\(^\text{[85]}\)、哈勃空间望远镜、以及威尔金森微波各向异性探测器 WMAP\(^\text{[86]}\)）的数据分析，大爆炸宇宙学取得了显著进展。如今，宇宙学家们已经相当精确地测定了大爆炸模型中的多个关键参数，并意外地发现——宇宙的膨胀似乎正在加速。\(^\text{[87][88]}\)
\subsection{观测证据}
大爆炸模型为广泛的观测现象提供了一个全面的解释，包括轻元素的丰度、宇宙微波背景辐射、大尺度结构以及哈勃定律。\(^\text{[90]}\)最早且最直接支持该理论的观测证据包括：根据哈勃定律所揭示的宇宙膨胀（由星系红移表明）、宇宙微波背景辐射的发现与测量、以及由“大爆炸核合成”（BBN）产生的轻元素的相对丰度。更近期的证据还包括星系的形成与演化观测，以及宇宙大尺度结构的分布。\(^\text{[91]}\)这些证据有时被称为“大爆炸模型的四大支柱”。\(^\text{[92][93]}\)

现代精确的大爆炸模型依赖于一些尚未在地球实验室中观测到、或尚未被纳入粒子物理标准模型（Standard Model）的奇异物理现象。在这些特征中，暗物质（dark matter）目前是实验研究最为活跃的领域。\(^\text{[94]}\)然而，冷暗物质模型仍存在“尖核晕问题”（cuspy halo problem）\(^\text{[95]}\)与“矮星系问题”（dwarf galaxy problem）\(^\text{[96]}\)。暗能量（dark energy）同样是科学家高度关注的领域，但目前尚不清楚能否实现暗能量的直接探测。\(^\text{[97]}\)“暴胀”（inflation）与“重子生成”（baryogenesis）仍属于大爆炸模型中较具推测性的部分。目前科学界仍在寻求可行且可定量化的解释，这些问题被视为物理学中的未解之谜。
\subsubsection{哈勃定律与宇宙膨胀}
\begin{figure}[ht]
\centering
\includegraphics[width=8cm]{./figures/705d2a074dce6061.png}
\caption{} \label{fig_BBang_8}
\end{figure}
对遥远星系与类星体的观测显示，这些天体的光谱呈现红移现象：它们发出的光被移向更长的波长。通过对目标进行频谱分析，可以将其发射或吸收线的光谱模式，与化学元素原子与光相互作用所产生的特征线相匹配。这些红移是各向同性的，在所有方向上观测到的天体中分布均匀。如果将红移解释为多普勒效应（Doppler shift），则可以计算出天体的退行速度（recessional velocity）。对于某些星系，还可以通过“宇宙距离梯”（cosmic distance ladder）估算它们的距离。当退行速度与距离绘制在坐标图上时，会出现一个线性关系，即著名的哈勃定律（Hubble’s law）：[63]$v = H_0 D$其中：  
\begin{itemize}
\item \( v \) 为星系或其他遥远天体的退行速度；  
\item \( D \) 为该天体的真实距离（proper distance）；  
\item \( H_0 \) 为哈勃常数（Hubble’s constant），根据 WMAP 的测量为 \( H_0 = 70.4^{+1.3}_{-1.4} \, \text{km/s/Mpc} \)。\(^\text{[40]}\)
\end{itemize}
哈勃定律表明宇宙在各处均匀膨胀。这种宇宙膨胀最早由弗里德曼（Alexander Friedmann）在 1922 年\(^\text{[61]}\)和勒梅特（Georges Lemaître）在 1927 年\(^\text{[65]}\)依据广义相对论所预言，远早于哈勃在 1929 年的分析与观测。这一发现成为大爆炸模型（由弗里德曼、勒梅特、罗伯逊与沃克发展）的基石。

该理论要求以下关系在任意时刻均成立：$v = H D$其中 \( D \) 为真实距离（proper distance），\( v \) 为退行速度。随着宇宙的膨胀，\( v \)、\( H \)、\( D \) 均随时间变化。因此，用 \( H_0 \) 表示“当今的哈勃常数”。在远小于可观测宇宙尺度的距离上，哈勃红移可以视为由退行速度 \( v \) 引起的多普勒位移。而在与可观测宇宙相当的尺度上，红移的成因更为复杂，但其作为运动学多普勒效应的解释仍然是最自然的。\(^\text{[98]}\)

在测定哈勃常数的过程中出现了尚未解释的不一致现象，被称为“哈勃张力”（Hubble tension）。基于宇宙微波背景（CMB）观测的测量结果表明哈勃常数值较低，而通过“宇宙距离梯”方法测得的结果则较高。\(^\text{[99]}\)
\subsubsection{宇宙微波背景辐射}
\begin{figure}[ht]
\centering
\includegraphics[width=8cm]{./figures/59e7fba922848e68.png}
\caption{由COBE 卫星上的FIRAS 仪器（Far Infrared Absolute Spectrophotometer）测得的宇宙微波背景辐射（CMB）光谱，是自然界中迄今为止测量最精确的黑体辐射谱。\(^\text{[100]}\)在该图中，实验数据点及其误差棒几乎被理论曲线完全覆盖。} \label{fig_BBang_9}
\end{figure}
1964 年，阿诺·彭齐亚斯（Arno Penzias）与罗伯特·威尔逊（Robert Wilson）偶然发现了宇宙背景辐射（cosmic background radiation）——  
一种在微波波段中全方位存在的信号。\(^\text{[79]}\)这一发现为阿尔弗（Ralph Alpher）、赫尔曼（Robert Herman）与伽莫夫（George Gamow）在 1950 年左右提出的“大爆炸预测”提供了有力的实验证实。到 1970 年代，科学家们发现这种辐射在各个方向上与黑体辐射谱（blackbody spectrum）近似一致；该谱线因宇宙膨胀而红移，如今对应的温度约为 2.725 K。这一结果使证据天平倾向于大爆炸模型，彭齐亚斯与威尔逊因此获得了 1978 年诺贝尔物理学奖。

与宇宙微波背景辐射（CMB）对应的“最后散射面”（surface of last scattering）出现在复合纪（recombination epoch）之后不久——即中性氢（neutral hydrogen）开始稳定形成的时期。在此之前，宇宙由炽热而致密的光子–重子等离子体（photon–baryon plasma）构成，光子不断被自由带电粒子散射，难以长距离传播。当宇宙年龄约为 \(372 \pm 14\) 千年（kyr）时，\(^\text{[101]}\)光子的平均自由程（mean free path）变得足够长，可以穿越宇宙直至今日，此时宇宙变得透明。
\begin{figure}[ht]
\centering
\includegraphics[width=8cm]{./figures/7f49c214e3e5245f.png}
\caption{2012 年公布的WMAP 九年宇宙微波背景辐射图像（9-year WMAP image of the cosmic microwave background radiation）。\(^\text{[102][103]}\)该辐射在各个方向上几乎完全均匀，其各向同性精确到约 十万分之一（1/100,000）。\(^\text{[104]}\)} \label{fig_BBang_10}
\end{figure}
1989 年，美国国家航空航天局（NASA）发射了COBE（宇宙背景探测者，Cosmic Background Explorer）卫星，并取得了两项重大突破：在 1990 年，COBE 的高精度光谱测量表明，宇宙微波背景辐射（CMB）的频谱几乎是一个完美的黑体谱，其偏差小于 \(10^{-4}\)；同时测得残余温度为 2.726 K（更近的测量结果略微修正为 2.7255 K）。随后在 1992 年，COBE 进一步观测到 CMB 温度在全天空中存在微小波动（各向异性，anisotropies），其幅度约为 \(10^{-5}\) 量级\(^\text{[85]}\)约翰·C·马瑟（John C. Mather）与乔治·斯穆特（George F. Smoot）因在这些研究中的领导作用，获得了2006 年诺贝尔物理学奖。

在接下来的十年间，大量地面与气球观测实验继续研究 CMB 的各向异性。在 2000–2001 年间，多个实验（尤其是BOOMERanG）通过测量各向异性的典型角尺度（即在天空中的表观尺寸），发现宇宙在空间几何上几乎是平坦的。\(^\text{[105][106][107]}\)

2003 年初，威尔金森微波各向异性探测器（WMAP, Wilkinson Microwave Anisotropy Probe）公布了首批观测结果，提供了当时最精确的一组宇宙学参数测量值。这些结果排除了若干特定的宇宙暴胀模型，但总体上仍与暴胀理论（inflation theory）相一致。\(^\text{[86]}\)随后在 2009 年 5 月，普朗克卫星（Planck space probe）成功发射。与此同时，其他地面与气球平台的宇宙微波背景实验仍在持续进行中。
\subsubsection{原初元素丰度}
\begin{figure}[ht]
\centering
\includegraphics[width=8cm]{./figures/82719a075d7d217f.png}
\caption{大爆炸核合成过程中轻元素丰度的时间演化} \label{fig_BBang_11}
\end{figure}
根据大爆炸模型（Big Bang models），可以计算出宇宙中氦-4（\(^4\mathrm{He}\)）、氦-3（\(^3\mathrm{He}\)）、氘（\(^2\mathrm{H}\)）和锂-7（\(^7\mathrm{Li}\)）等同位素相对于普通氢（H）的预期丰度比。\(^\text{[37]}\)这些相对丰度仅依赖于一个参数——光子与重子（baryon）的比值。该比值还可通过宇宙微波背景（CMB）温度涨落的细节结构独立计算得到。模型预测的质量比约为：\(^4\mathrm{He}\!:\!H \approx 0.25\)，\(^2\mathrm{H}\!:\!H \approx 10^{-3}\)，\(^3\mathrm{He}\!:\!H \approx 10^{-4}\)，以及 \(^7\mathrm{Li}\!:\!H \approx 10^{-9}\)。[37]

观测得到的各元素丰度大体上与单一重子-光子比值下的预测结果一致。氘（\(^2\mathrm{H}\)）的符合度极好，氦-4（\(^4\mathrm{He}\)）略有偏差但总体接近，而锂-7（\(^7\mathrm{Li}\)）则偏差约一倍——这一异常被称为宇宙学锂问题（cosmological lithium problem）。在后两种情况下，仍存在显著的系统性不确定性。尽管如此，与大爆炸核合成（BBN）预测的总体一致性仍然是支持大爆炸理论的强有力证据，因为该理论是目前唯一能够解释轻元素相对丰度的模型，而且几乎不可能通过调节模型参数使氦的产量远离 20–30\% 这一范围。\(^\text{[108]}\)事实上，除了大爆炸之外，没有其他明显理由能够解释为何在恒星形成之前的年轻宇宙中（通过观测被认为不含恒星核合成产物的物质所确定）氦应多于氘，而氘又多于氦-3，且这些比例在不同观测中几乎保持恒定。\(^\text{[109]: 182–185 }\)
\subsubsection{星系的演化与分布}
对星系与类星体（quasar）的形态和分布的详细观测结果与当前大爆炸模型的预测一致。观测与理论的综合研究表明，最早的类星体与星系在大爆炸后一十亿年内形成，\(^\text{[110]}\)此后逐渐出现更大的结构，如星系团与超星系团。\(^\text{[111]}\)

恒星族群不断老化与演化，因此距离较远的星系（即观测到的早期宇宙状态）与近邻星系（较晚期状态）在外观上差异显著。此外，较晚形成的星系与同样距离但在大爆炸后不久形成的星系相比，其结构也有显著不同。这些观测结果是反对稳态宇宙模型（steady-state model）的重要证据。对恒星形成率、星系与类星体分布、以及大尺度结构的观测与大爆炸模型下的宇宙结构形成模拟结果高度吻合，并持续完善该理论的细节。\(^\text{[111][112]}\)
\subsubsection{原初气体云}
\begin{figure}[ht]
\centering
\includegraphics[width=8cm]{./figures/a2de458976f8cfcb.png}
\caption{显微镜下的 BICEP2 望远镜焦平面 ——  用于探测宇宙微波背景辐射（CMB）中的偏振信号。\(^\text{[113][114][115][116]}\)} \label{fig_BBang_12}
\end{figure}
2011年，天文学家通过分析遥远类星体光谱中的吸收线，发现了他们认为是原始气体（primordial gas）的原初气体云。在此发现之前，所有被观测到的天体都含有由恒星核合成形成的重元素。然而，尽管观测仪器对碳、氧、硅极为敏感，在这两团气体云中却未检测到这三种元素的存在。\(^\text{[117][118]}\)由于这些气体云未检测到可观的重元素含量，它们很可能在大爆炸后的最初几分钟内形成，即在大爆炸核合成（BBN）阶段产生的。
\subsubsection{其他证据}
根据哈勃膨胀率和宇宙微波背景（CMB）推算出的宇宙年龄，如今与其他独立的年龄估算方法相一致。这些方法包括：—— 通过将恒星演化理论应用于球状星团，以及对第二星族（Population II）恒星进行放射性测年（radiometric dating）。\(^\text{[119]}\)此外，它还与基于Ia型超新星（Type Ia supernovae）观测到的膨胀速率以及宇宙微波背景温度涨落测量所得的年龄估算结果相符。\(^\text{[120]}\)

不同观测方法得到的宇宙年龄一致，进一步支持了ΛCDM 模型（Lambda-Cold Dark Matter model），因为该模型用于将部分观测量与宇宙年龄联系起来，而所有独立的推算结果都互相吻合。不过，某些早期宇宙天体（尤其是类星体 APM 08279+5255）的观测仍引发了一些质疑——这些天体在 ΛCDM 模型下是否真的有足够时间在如此早期形成。\(^\text{[121][122]}\)

模型还预测：宇宙微波背景辐射温度在过去更高，这一点已通过对高红移气体云中极低温吸收线的观测得到实验证实。\(^\text{[123]}\)该预测还意味着：星系团中 Sunyaev–Zel'dovich 效应的振幅与红移无直接依赖关系。观测表明这一点大体成立，但由于星系团性质会随宇宙时间演化而改变，精确测量这一效应仍然困难。\(^\text{[124][125]}\)
\subsubsection{未来观测}
未来的引力波观测台（gravitational-wave observatories）有望探测到原初引力波（primordial gravitational waves）—— 即来自早期宇宙、可追溯至大爆炸后不到一秒的遗迹信号。\(^\text{[126][127]}\)
\subsection{物理学中的问题与相关议题}
与所有理论一样，大爆炸模型的发展也带来了一系列谜团与难题。其中一些已被解决，而另一些仍未有定论。为了解决这些问题而提出的假说，又往往引出了新的谜题。例如，视界问题（horizon problem）、磁单极问题（magnetic monopole problem）以及平坦性问题（flatness problem），通常可借助暴胀理论（inflation theory）来解释，但暴胀宇宙的细节仍未厘清，甚至有部分理论奠基者声称该理论已被证伪。\(^\text{[128][129][130][131]}\)以下列出的是目前宇宙学家与天体物理学家仍在深入研究的与“大爆炸”概念相关的主要谜题。
\subsubsection{重子不对称性}
目前尚不清楚为什么宇宙中物质多于反物质。\(^\text{[34]}\)一般假设在早期炽热的宇宙中，系统处于统计平衡，含有相等数量的重子与反重子。当时物质与反物质的总量远高于今天，仅存在约百亿分之一的微小不对称。随着物质与反物质相互湮灭，只留下这极少数剩余的物质。当今的观测表明，即使在宇宙最遥远的区域，也几乎完全由普通物质构成，而反物质极为稀少。\(^\text{[132]}\)

若物质与反物质完全对称，湮灭将只留下光子，而几乎不残留任何物质——显然与观测不符。为解释这种不对称，物理学家提出了重子生成（baryogenesis）过程。要发生重子生成，必须满足萨哈罗夫条件（Sakharov conditions），即：重子数不守恒；存在 C 对称性与 CP 对称性破缺；宇宙偏离热力学平衡。\(^\text{[133][134]}\)这些条件在标准模型中均可出现，但其效应强度不足以解释当今观测到的重子不对称性。
\subsubsection{暗能量}
对 Ia 型超新星红移—亮度关系的测量表明，自宇宙约为当前年龄的一半起，宇宙膨胀速度便开始加速。为了说明这种加速，宇宙学模型假设宇宙能量的大部分由一种具有巨大负压（negative pressure）的成分组成，称为暗能量（dark energy）。\(^\text{[11]}\)

尽管暗能量的存在仍属推测，但它可解决多个难题。宇宙微波背景测量显示宇宙空间几乎是平坦的（spatially flat），因此根据广义相对论，宇宙的总质量/能量密度必须接近临界密度。然而，通过引力聚集测得的物质密度仅为临界密度的约 30\%。\(^\text{[11]}\)由于暗能量被认为不会像常规物质那样聚集，它被视为“缺失能量密度”的最佳解释。  

暗能量还能解释宇宙整体曲率的两个几何测量指标：引力透镜（gravitational lenses）的频率；\(^\text{[135]}\)重子声学振荡（baryon acoustic oscillations, BAO）的特征分布，即宇宙学中的“宇宙标尺”。\(^\text{[136][137]}\)

负压被认为是真空能（vacuum energy）的一种性质，但暗能量的确切本质与存在形式仍是大爆炸理论中最大的未解之谜之一。WMAP 团队在 2008 年的结果表明，宇宙由73\% 的暗能量、23\% 的暗物质、4.6\% 的普通物质，以及不足 1\% 的中微子组成。\(^\text{[40]}\)理论上，物质能量密度会随宇宙膨胀而降低，但暗能量密度保持恒定（或几乎恒定）。因此，在过去物质占宇宙总能量的比例较大，但在遥远的未来，随着宇宙继续膨胀，暗能量将变得更加主导。

理论物理学家用多种相互竞争的理论来解释暗能量成分，包括爱因斯坦的宇宙学常数（cosmological constant）、第五元素假说（quintessence）以及其他修正引力理论。\(^\text{[138]}\)由此衍生的宇宙学常数问题（cosmological constant problem），有时被称为“物理学中最令人尴尬的问题”，源自于观测到的暗能量密度与按普朗克单位（Planck units）所作的理论预估值之间存在极大的差异。\(^\text{[139]}\)
\subsubsection{暗物质}
\begin{figure}[ht]
\centering
\includegraphics[width=8cm]{./figures/dec3cdee679f28f2.png}
\caption{图表显示了宇宙不同成分所占的比例——其中约 95\% 为暗物质与暗能量。} \label{fig_BBang_13}
\end{figure}
在 1970 至 1980 年代的观测中，天文学家发现宇宙中可见物质（visible matter）的量不足以解释星系内部与星系之间观测到的引力作用强度。这引出了一个假设：宇宙中多达90\% 的物质是“暗物质（dark matter）”，它既不发光，也不与普通重子物质发生电磁相互作用。  

此外，若假定宇宙主要由普通物质构成，理论预测的结果会与实际观测严重不符。  
特别是，若无暗物质存在，宇宙的结构应当更加均匀且富含氘，而现实中我们观测到的宇宙更加团块化（lumpy），且氘的丰度显著更低。尽管暗物质一直存在争议，但多项观测结果都暗示它的存在：包括宇宙微波背景各向异性（CMB anisotropies）、星系旋转曲线问题（galaxy rotation problem）、星系团速度弥散（velocity dispersions）、大尺度结构分布（large-scale structure distributions）、 引力透镜观测（gravitational lensing）以及星系团的 X 射线测量等。\(^\text{[140]}\)

暗物质的间接证据主要来自它对其他物质的引力影响，因为迄今为止尚无任何暗物质粒子在实验室中被直接探测到。粒子物理学家提出了多种暗物质候选粒子，多个国际实验（如 LUX、XENON、PandaX）正在尝试直接探测它们。\(^\text{[141]}\)

然而，当前主流的冷暗物质模型（Cold Dark Matter, CDM）仍面临若干尚未解决的问题，其中包括矮星系问题（dwarf galaxy problem）\(^\text{[96]}\)以及尖峰晕问题（cuspy halo problem）\(^\text{[95]}\)。因此，也有学者提出不依赖暗物质的替代理论，通过修正牛顿或爱因斯坦的引力定律（Modified Gravity）来解释观测现象。但迄今为止，尚无任何替代理论能像冷暗物质模型那样成功解释所有观测结果。\(^\text{[142]}\)
\subsubsection{视界问题}
视界问题源于这样一个前提：信息传播速度不能超过光速。在一个具有有限年龄的宇宙中，这一限制导致存在一个粒子视界（particle horizon），它决定了任意两片空间区域之间是否能通过因果联系相互影响。\(^\text{[143]}\)然而，宇宙微波背景的高度各向同性（isotropy）在此框架下难以解释：若宇宙自始至终都由辐射或物质主导，则在最后散射面时，粒子视界仅对应约2°的天空角尺度。因此，超出这一角度的天空区域之间不应具有相同温度。\(^\text{[109]: 191 }\)

这一表面矛盾可由暴胀理论（inflation theory）解释。根据该理论，在极早期（重子生成之前），宇宙被一个均匀且各向同性的标量能量场（scalar energy field）所主导。在暴胀阶段，宇宙经历指数级膨胀（exponential expansion），粒子视界的扩张速度远超此前假设。因此，如今位于可观测宇宙两端的区域，在暴胀之前实际上处于同一因果区域之内。这解释了 CMB 的温度一致性：这些区域在暴胀前就已热平衡。\(^\text{[28]: 180 }\)

根据海森堡不确定性原理（Heisenberg's uncertainty principle），暴胀阶段会产生量子热涨落（quantum thermal fluctuations），这些微小涨落被放大至宇宙尺度，成为后来宇宙中一切结构的“种子”。\(^\text{[109]: 207 }\)暴胀理论预测，这些原初涨落应当是近似尺度不变（scale invariant）且服从高斯分布（Gaussian）的，这一点已通过 CMB 的高精度观测得到验证。\(^\text{[86]: sec 6 }\)

此外，一个与经典视界问题相关的难题在于：多数标准宇宙暴胀模型预言，暴胀结束早于电弱对称性破缺（electroweak symmetry breaking）发生的时间。因此，暴胀不应能够消除因果上分离区域中电弱真空的大尺度不连续性（large-scale discontinuities）。\(^\text{[144]}\)
\subsubsection{磁单极问题}
在 1970 年代后期，物理学家提出了磁单极问题（magnetic monopole problem）。根据大统一理论（Grand Unified Theories, GUTs），空间中应存在拓扑缺陷（topological defects），表现为磁单极（magnetic monopoles）。在炽热的早期宇宙中，这类粒子应被高效产生，导致其数密度远高于观测允许的水平，而事实上我们未探测到任何磁单极的存在。这一问题由宇宙暴胀理论得以解决：暴胀将所有点状缺陷从可观测宇宙中“稀释”出去，就如同它令宇宙的几何趋于平坦一样。\(^\text{[143]}\)
\subsubsection{平坦性问题}
\begin{figure}[ht]
\centering
\includegraphics[width=8cm]{./figures/c50b8969928a0bb4.png}
\caption{宇宙的整体几何结构由宇宙学参数 \( \Omega \)（欧米伽参数）是小于、等于还是大于 1 来决定。从上到下依次显示：正曲率的闭合宇宙（closed universe）、负曲率的双曲宇宙（hyperbolic universe），以及零曲率的平坦宇宙（flat universe）。} \label{fig_BBang_14}
\end{figure}
平坦性问题（Flatness problem，又称“老龄问题”Oldness problem）是与弗里德曼–勒梅特–罗伯逊–沃克（FLRW）宇宙模型相关的一个观测难题。\(^\text{[143]}\)宇宙的空间曲率可能为正、负或零，这取决于其总能量密度（total energy density）。当能量密度低于临界密度时，曲率为负；高于临界密度时，曲率为正；恰等于临界密度时，曲率为零，此时空间被称为平坦（flat）。观测结果表明，宇宙与平坦状态高度一致。\(^\text{[145][146]}\)

问题在于：任何微小偏离临界密度的差异，都会随着时间而放大，然而今天的宇宙仍然极其接近平坦。[notes 2]若以普朗克时间（Planck time）约 \(10^{-43}\) 秒作为自然时间尺度，\(^\text{[33]}\)那么宇宙在经历数十亿年后既未坍缩为“大挤压（Big Crunch）”，也未热寂（heat death），显然需要额外的解释。例如，即便在相对较晚的时期——大约几分钟时（核合成阶段），宇宙的密度也必须与临界密度相差不超过 \(10^{14}\) 分之一，否则宇宙便不会以如今的形式存在。\(^\text{[147]}\)
\subsection{误解}
除了人们对宇宙膨胀性质的普遍混淆外，大爆炸模型本身也常被误解。

一个常见的误解是认为大爆炸模型**完全解释了宇宙的起源**。  
事实上，大爆炸模型并未描述**能量、时间与空间是如何被“创造”**的，  
而是描述了当今宇宙如何从一种**超高密度与高温的初始状态**中演化而来。[148]

另一个常见误解涉及**哈勃定律（Hubble’s law）**中退行速度的理解。  
这些速度并非相对论意义上的速度  
（例如，它们并不对应于四维速度的空间分量）。  
因此，根据哈勃定律，距离超过哈勃长度的星系**以超过光速的速率退行**，  
并不意味着这些星系正在以超光速“运动”，  
这种“速度”并不等同于超光速旅行。[149]

启示（Implications）

按照当前的科学认识，对宇宙未来的推断  
仅能在有限的时间范围内进行——  
尽管这个时间尺度远远长于宇宙目前的年龄，  
但超出此范围后的推测将变得越来越具有猜测性。  
同样地，目前对于宇宙起源的真正理解  
仍然只能停留在**理论假设与推测**的层面。[150]
```




